\section{方法}
\label{sec:method}

\subsection{可迁移性边界与粒度失配}

若解码期干预模块的输入或输出具有宿主特异性，则难以称为严格通用。早期校准方法读取模型原生隐状态并投影到宿主词表几何中，此类可学习映射按构造即与宿主耦合，无法在无重训情况下迁移到异构 MLLM。

要实现真正的热插拔，干预须运行于标准化的外部语义空间。设计外部控制器又必须直面幻觉的根本动因之一：\textbf{跨模态粒度失配}。MLLM 将图像编码为连续、整体化的特征图，而自回归生成却输出离散、常呈碎片化的子词 token（例如 \texttt{["gi", "\#\#raf", "\#\#fe"]}）。在连续视觉实体与无语义碎片 token 之间建立可靠物理接地极为困难。粒度崩塌使模型陷入\emph{开环解码}：视觉校验失效，语言先验惯性压倒视觉约束。CHORD 旨在弥合该失配并闭合解码回路。

\subsection{CHORD 框架概述}

CHORD（Calibrating Hallucinations via Object-Resonant Decoding）为免训练、即插即用解码框架。它不覆盖全词表分布，仅作用于宿主在步 $t$ 的活跃解码前沿
\begin{equation}
  \mathcal{F}_t = \mathrm{Top}\text{-}K(L_t),
\end{equation}
其中 $L_t$ 为原始 logit。流程含四阶段：（1）分词器桥接；（2）连续共振监控；（3）假设驱动共振增量；（4）logit 校准。

生成前，CHORD 对输入图像 $I$ 做一次离线提取：冻结的外部视觉编码器构建稠密局部区域特征库 $\mathcal{Z}_R = \{z_{R_1}, z_{R_2}, \dots, z_{R_N}\}$，作为后续共振核验的静态``物理事实''缓存。

\subsection{无参数分词器桥接}

为克服粒度失配，CHORD 不直接对原始 token ID 评分。设 $B_m$ 为与宿主 $m$ 关联的\textbf{无参数}分词器桥接。对每个候选 $c_k \in \mathcal{F}_t$，桥接利用分词器前缀树做确定性前瞻，将 token 提升为语义完整的文本片段 $A_k$，得到候选片段集 $\mathcal{S}_t = \{A_1,\dots,A_K\}$。桥接不含可学习参数，是将模型特异的碎片化提议转为跨模型、可物理接地的语义命题的关键。

\subsection{假设驱动的共振增量}

CHORD 将解码建模为物理共振过程：将前缀 $p$ 与候选片段 $A_k$ 拼接视为指向区域库 $\mathcal{Z}_R$ 的``语义声纳''。为严格避免\emph{空间错配}（前缀与候选对齐到无关区域），共振增量锚定在\textbf{同一}物理区域上。首先为联合假设选取最相关区域：
\begin{equation}
  r^* = \arg\max_{r \in \mathcal{Z}_R} \cos\bigl(E_{\text{text}}(p \oplus A_k), z_r\bigr),
\end{equation}
其中 $E_{\text{text}}$ 为冻结公开文本编码器。再定义\textbf{共振增量 $\Delta S$}，刻画加入 $A_k$ 的纯物理效应：
\begin{equation}
  \Delta S(A_k) = \cos\bigl(E_{\text{text}}(p \oplus A_k), z_{r^*}\bigr) - \cos\bigl(E_{\text{text}}(p), z_{r^*}\bigr).
\end{equation}

该式起到\textbf{无启发式}判别作用：（1）$\Delta S>0$ 表示视觉支持或正确空间关系；（2）$\Delta S<0$ 表示与图像矛盾（如虚构物体）；（3）$\Delta S\approx 0$ 对应抽象词或语法词，CHORD 可安全忽略而无需脆弱 POS 标注。

\subsection{克服过度自信：连续监控与系统优化}

朴素熵门控难以捕获\textbf{过度自信幻觉}（低熵仍错）。CHORD 弃用纯熵门控，采用\textbf{连续 Top-$1$ 监控与全 $K$ 回退}，并结合 $E_{\text{text}}$ 的 KV 缓存。前缀 $p$ 持续缓存；$p\oplus A_k$ 往往只需再前向 1--2 个 token。每步对宿主最自信 token 计算 $\Delta S(A_{\text{top1}})$：若 $\Delta S(A_{\text{top1}})\ge 0$ 则直接接受；若 $<0$ 则触发对全部 $K$ 个候选的完整共振评估以寻找接地替代。该策略在抑制过度自信的同时控制计算开销。

\subsection{Logit 校准}

触发完整干预时，CHORD 将共振证据注入宿主 logit。与片段 $A_k$ 对应的候选 token $c_k$ 的调整 logit 为
\begin{equation}
  \tilde{L}_t(c_k) = L_t(c_k) + \alpha \cdot \Delta S(A_k),
\end{equation}
$\alpha$ 为强度超参；再经 softmax 归一化。有界重分配使轨迹向视觉事实偏转，同时尽量不破坏宿主语言结构且无需更新参数。
